\section{结论与局限}

本报告在纽约签到驱动的派单仿真上完成了收入对比、帕累托前沿和两项拓展实验，结论如下。(1) 长期价值的作用取决于运力是否紧张。Q-KM 在 40/60 辆车、各批次间隔下显著优于 KM 即时收入与最近贪心，增收来自“白天保持空间分布、晚高峰兑现运力”这一跨时段取舍；100 辆车时不再有优势。批次间隔一旦超过订单耐心，收入会骤降。(2) 派单的双目标前沿以非支撑解为主，加权和方法在原理上只能覆盖前沿的一小部分。以精确前沿为真值，NSGA-II 优于 NSGA-III；用 KM 解初始化种群后，近似质量进一步提高，并且统计显著。(3) Q-KM 在大 $\beta$ 下的失效主要来自 max 自举的过估计，期望型的策略评估目标可以消除这一问题。深度价值网络提供了对 $\beta$ 的鲁棒性和连续的价值面，前提是加入空间平滑正则，但在本数据规模下没有带来更高的收入峰值。进一步的乐观度谱实验表明，崩溃并不取决于高估的幅度，而取决于高估在“等待”与“出发”之间是否对称；长期价值派单需要的是动作间相对价值无偏。(4) \HINDSIGHTCONC (5) 单批次前沿会把收入–接驾的权衡代价高估数倍：由于批次间的订单替代和运力回收，较小的接驾惩罚几乎不损失全天收入，多目标派单应在系统层面评估。

\textbf{局限。} 订单由签到数据构造，与真实打车需求存在偏差；仿真中空闲司机不会主动调度；价值函数与精确前沿都只针对单个批次或单一城市。今后可以把价值网络与空车调度联合优化，并把 Q-KM 的长期价值边权作为第三个目标，在全天层面研究“收入–接驾距离–应答率”的系统级前沿。
